\documentclass[10pt,a4paper]{amsart}
\usepackage[margin=19mm]{geometry}
\usepackage{amsmath,amssymb,mathtools}
\usepackage[scr=rsfs,cal=euler]{mathalfa}
\usepackage{xfrac}
\usepackage[unicode,hidelinks,bookmarks=false]{hyperref}
\newcommand{\ie}{i.e.\ }
\newcommand{\cf}{cf.\ }
\newcommand{\resp}{respectively\ }
\newcommand{\Subsection}{\subsection}
\providecommand{\nobreakdash}{\nobreak\hbox{-}}
\setlength{\emergencystretch}{2em}
\multlinegap0pt
\usepackage{amssymb}
\usepackage{tikz}
\usepackage{enumerate}
\newtheorem{theorem}{Theorem}
\renewcommand\geq{\geqslant}
\title{Bounds for characteristic numbers of conic--line arrangements in the plane}
\author{Rita Pardini, Piotr Pokora}
\date{Source: arXiv:2312.12950v2 (CC BY 4.0)}
\begin{document}
\maketitle

\noindent\emph{Source and licence.} Bounds for characteristic numbers of conic--line arrangements in the plane. Version arXiv:2312.12950v2 (CC BY 4.0), distributed by arXiv under the Creative Commons Attribution 4.0 International licence, http://creativecommons.org/licenses/by/4.0/, as recorded in the arXiv licence metadata for this exact version and shown on the abstract page https://arxiv.org/abs/2312.12950v2. Full licence notice: \texttt{licenses/arxiv-2312.12950-CC-BY-4.0.txt}.\par\smallskip

\noindent\textit{Editorial scope.} Counted unit: the article's theorem theorem (source0.tex lines 522--527). The article's complete original proof of it is delivered with it (source0.tex lines 528--561, one \texttt{proof} environment of 34 lines). No mathematics is written, repaired or re-derived: the statement and the proof are the article's own lines, in the article's own notation, and no retained line is substituted or re-flowed.  No further source-local context is needed: every cross-reference of every retained line resolves inside the two delivered spans.  Nothing was omitted from any retained span, so the package carries no omission record.  No retained line contains a citation, so the wrapper carries no bibliography and no bibliography entry.  No retained line contains a figure environment, an image-inclusion command or drawing code, so no image is delivered and the package declares no asset dependency.  Editorial formatting only, in addition: an AMS \texttt{amsart} preamble that restates the article's own declarations and macros used by the retained lines, namely the environments \texttt{theorem} on separate counters, together with the standard packages those lines need.  The retained environments print in this document as Theorem 1 in order of appearance, whereas the article prints the same items with the numbering of its own full document.  The complete native source of the article as distributed in the arXiv e-print for this exact version is bundled unchanged as \texttt{sources/151819/source0.tex}.
\begin{theorem}
Let $\mathcal{C}$ be an arrangement of $k\geq 3$ smooth conics in the plane such that $t_{k}=0$.
Then 
$$\gamma(\mathcal{C}) :=\frac{9-8k+3f_{1}-4f_{0}}{3-2k+f_{1}-f_{0}} \geq \frac{4k^2-12k+9}{2k^{2}-4k+3}.$$
Moreover, the lower bound is achieved when our arrangement has only double points as intersections.
\end{theorem}

\begin{proof}
Observe that if $\mathcal{C}$ admits only double intersection points, then
$$\gamma(\mathcal{C})= \frac{9-8k+2t_{2}}{3-2k+t_{2}} = \frac{4k^2-12k+9}{2k^{2}-4k+3}.$$
Now we want to show that the case when we have only double intersection points as singularities gives  a lower bound  for $\gamma(\mathcal{C})$. We wish  to show that
\begin{equation}
\frac{9-8k+3f_{1}-4f_{0}}{3-2k+f_{1}-f_{0}} \geq \frac{4k^2-12k+9}{2k^{2}-4k+3}.
\end{equation}
Observe that the above inequality is equivalent to
\begin{equation}
\label{main}
-8k^3+14k^2-6k + 2k^2f_{1}+(-4k^2+4k-3)f_{0} \geq 0.     
\end{equation}
Recall that the following combinatorial count holds
$$4\cdot \binom{k}{2} = 2(k^2-k) = \sum_{r\geq 2} \binom{r}{2}t_{r}.$$
Multiplying by $2$ the above formula we get
\begin{equation}
4k^2-4k  = \sum_{r \geq 2}r^{2}t_{r} - \sum_{r\geq 2} rt_{r}= f_{2} - f_{1}.
\end{equation}
and observe that
$$-8k^3+14k^2-6k = -2k(f_{2}-f_{1}) + \frac{3}{2}(f_{2}-f_{1}) = \bigg(-2k+\frac{3}{2}\bigg)f_{2} + \bigg(2k-\frac{3}{2}\bigg)f_{1}.$$
Plugging this into \eqref{main}, we obtain
$$\bigg(-2k + \frac{3}{2}\bigg)f_{2}+ \bigg(2k^2 + 2k-\frac{3}{2}\bigg)f_{1} +(-4k^2 + 4k - 3)f_{0} =$$
$$\sum_{r\geq 2}\bigg( \bigg(-2k+\frac{3}{2}\bigg)r^{2} + \bigg(2k^{2} + 2k - \frac{3}{2}\bigg)r + (-4k^{2} + 4k-3)\bigg)t_{r} \geq 0.$$
In order to finish the proof, we have to show that 
\begin{equation}
\label{check}
\bigg( \bigg(-2k+\frac{3}{2}\bigg)r^{2} + \bigg(2k^{2} + 2k - \frac{3}{2}\bigg)r + (-4k^{2} + 4k-3)\bigg) \geq 0
\end{equation}
for $r \in \{2, \ldots , k-1\}$ and suitably taken values of $k$, i.e., for $r \in \mathbb{N}_{\geq 2}$ the inequality must  hold with $k\geq r + 1$, since $t_{k}=0$. If we plug in $k = r + h$ with $h\geq 1$, then the left-hand side of \eqref{check} has the following form:
\begin{equation}
(r-2)\cdot(4h(h+r-1)-r+3).
\end{equation}
Since $r\geq 2$ and $h\geq 1$ we see that $(r-2)\cdot(4h(h+r-1)-r+3) \geq 0$, which completes our proof.
\end{proof}

\end{document}
